\indent ವಿಭಕ್ತಿ ಪ್ರತ್ಯಯಗಳನ್ನು ಎಂಟು ವಿಧಗಳಾಗಿ ವಿಂಗಡಿಸಬಹುದು. ಅವುಗಳು ಯಾವುವೆಂದರೆ—\par

\vspace{5pt plus 4pt minus 1pt}
\noindent{\renewcommand{\arraystretch}{1.24}\setlength{\tabcolsep}{4pt}\rowcolors{2}{kbvrow}{white}
\begin{tabularx}{\textwidth}{X X X}
{\kbvsecfont\color{kbvaccentdark}ವಿಭಕ್ತಿ} & {\kbvsecfont\color{kbvaccentdark}ವಿಭಕ್ತಿ ಪ್ರತ್ಯಯ} &  \\[1pt]
\arrayrulecolor{kbvline}\specialrule{0.5pt}{0pt}{2pt}\arrayrulecolor{black}
(೧) ಪ್ರಥಮಾ ವಿಭಕ್ತಿ & ... & ಉ \\
(೨) ದ್ವಿತೀಯಾ ವಿಭಕ್ತಿ & ... & ಅನ್ನು \\
(೩) ತೃತೀಯಾ ವಿಭಕ್ತಿ & ... & ಇಂದ \\
(೪) ಚತುರ್ಥಿ ವಿಭಕ್ತಿ & ... & ಇಗೆ, ಗೆ, ಕ್ಕೆ, ಅಕ್ಕೆ \\
(೫) ಪಂಚಮೀ ವಿಭಕ್ತಿ & ... & ದೆಸೆಯಿಂದ \\
(೬) ಷಷ್ಠಿ ವಿಭಕ್ತಿ & ... & ಎಂ, ಅ \\
(೭) ಸಪ್ತಮೀ ವಿಭಕ್ತಿ & ... & ಅಲ್ಲಿ \\
(೮) ಸಂಬೋಧನಾ ವಿಭಕ್ತಿ & ... & ಎ, ಏ—ಇತ್ಯಾದಿ \\
\end{tabularx}}
\vspace{5pt plus 4pt minus 1pt}

\kbvsechead{\textbf{ವಿಭಕ್ತಿ ರೂಪಗಳು}}

\vspace{4pt plus 4pt minus 1pt}\noindent{\kbvsecfont\fontsize{11pt}{14.5pt}\selectfont\color{kbvaccentdark} (೧)\hspace{0.45em}ಪ್ರಥಮಾ ವಿಭಕ್ತಿ:}\hspace{0.5em}ಪ್ರಥಮಾ ವಿಭಕ್ತಿ ಪ್ರತ್ಯಯ “ಉ” ಆಗಿದೆ. ಈ ವಿಭಕ್ತಿ ಪ್ರತ್ಯಯವು ಸಾಮಾನ್ಯವಾಗಿ ಕರ್ತೃಪದದ ಮೇಲೆ ಬರುವುದು.\par

\vspace{3pt plus 2pt}\noindent\hspace*{1.2em}\kbvlabel{ಉದಾ}\par\vspace{1pt}

\vspace{2.5pt plus 3pt minus 0.5pt}\noindent\hspace*{1.2em}\kbvnum{1.}\hspace{0.55em}ಸೂರ್ಯನು ಪೂರ್ವದಲ್ಲಿ ಉದಯಿಸಿದನು.\par

\indent (ಇಲ್ಲಿ ಸೂರ್ಯ- + ನ್ + ಉ = ಸೂರ್ಯನು)\par

\vspace{2.5pt plus 3pt minus 0.5pt}\noindent\hspace*{1.2em}\kbvnum{2.}\hspace{0.55em}ಎಮ್ಮೆಯು ಹಾಲನ್ನು ಕೊಡುತ್ತದೆ.\par

\begin{kbvrule}
\noindent\kbvbullet\hspace{0.5em}ಹಲವು ಸಂದರ್ಭಗಳಲ್ಲಿ ಪ್ರಥಮಾ ವಿಭಕ್ತಿಯ ಏಕವಚನ ರೂಪವನ್ನು ನಾಮಪ್ರಕೃತಿ ರೂಪದಲ್ಲೇ ಹೇಳುವುದು ವಾಡಿಕೆ. ಅಂದರೆ ನಾಮಪ್ರಕೃತಿಯ ರೂಪವೇ ಪ್ರಥಮಾ ವಿಭಕ್ತಿಯ ರೂಪವಾಗಿರುವುದು.
\end{kbvrule}

\vspace{3pt plus 2pt}\noindent\hspace*{1.2em}\kbvlabel{ಉದಾ}\par\vspace{1pt}

\vspace{2.5pt plus 3pt minus 0.5pt}\noindent\hspace*{1.2em}\kbvnum{1.}\hspace{0.55em}ಚಂದ್ರ ಆಕಾಶದಲ್ಲಿ ಹೊಳೆಯುತ್ತಾನೆ.\par

\vspace{2.5pt plus 3pt minus 0.5pt}\noindent\hspace*{1.2em}\kbvnum{2.}\hspace{0.55em}ಗಂಧಕ ನೀರಲ್ಲಿ ಕರಗುವುದಿಲ್ಲ.\par

\vspace{4pt plus 4pt minus 1pt}\noindent{\kbvsecfont\fontsize{11pt}{14.5pt}\selectfont\color{kbvaccentdark} (೨)\hspace{0.45em}ದ್ವಿತೀಯಾ ವಿಭಕ್ತಿ:}\hspace{0.5em}ದ್ವಿತೀಯಾ ವಿಭಕ್ತಿ ಪ್ರತ್ಯಯ “ಅನ್ನು” ಆಗಿದೆ. ಈ ವಿಭಕ್ತಿ ಪ್ರತ್ಯಯವು ಸಾಮಾನ್ಯವಾಗಿ ಕರ್ಮಪದದ ಮೇಲೆ ಬರುವುದು.\par
